\documentclass[10pt]{beamer}

\usetheme{metropolis}
\usepackage{amsmath, amssymb, bm}
\usepackage{booktabs}
\usepackage{hyperref}

\title{Relative Standard Error in OLS for A/B Tests}
\author{David Masip, DS Manager @ DH Pricing}
\date{}

\begin{document}

\maketitle

% ============================================================
\section{Motivation}
% ============================================================

\begin{frame}{Setup: from absolute to relative effect}
We run an A/B test and fit an OLS regression:
\[
Y_i = \alpha + \tau \cdot T_i + \mathbf{X}_i^\top \bm{\beta} + \epsilon_i .
\]
\begin{itemize}
  \item $\tau$ = \alert{absolute} average treatment effect (ATE).
  \item Stakeholders usually want a \alert{relative} lift:
  \[
    L \;=\; \frac{\tau}{\bar{Y}_{\text{control}}}.
  \]
  \item We need a point estimate \emph{and} a standard error for $L$.
  \item Without covariates ($\mathbf{X}_i$ omitted), this OLS is equivalent to the classical two-sample $t$-test for difference in means.
\end{itemize}
\end{frame}

\begin{frame}{``divide the SE by the control mean'' is a common shortcut}
Tempting shortcut:
\[
\hat{L} = \frac{\hat{\tau}}{\bar{y}_{\text{control}}}, \qquad
\widehat{\mathrm{SE}}(L) \stackrel{?}{=} \frac{\mathrm{SE}(\hat{\tau})}{\bar{y}_{\text{control}}}.
\]
Why this is \alert{wrong in general}:
\begin{itemize}
  \item The shortcut is the identity
  \[
  \mathrm{SE}\!\left(\frac{A}{B}\right) \;\stackrel{?}{=}\; \frac{\mathrm{SE}(A)}{B},
  \]
  which only holds when $B$ is a \alert{known constant}.
  \item But $\bar y_{\text{control}}$ is an estimate: it has its own standard error.
  \item Confidence intervals built this way are miscalibrated (too narrow, overly confident).
\end{itemize}
\end{frame}

\begin{frame}{When does it accidentally work?}
The naive formula is approximately correct only when:
\begin{itemize}
  \item $\mathrm{SE}(\bar{Y}_{\text{control}}) / \bar{Y}_{\text{control}} \ll \mathrm{SE}(\hat\tau)/|\hat\tau|$ (control mean is very precisely estimated relative to the effect), and
  \item $\mathrm{Cov}(\hat\tau, \bar{Y}_{\text{control}}) \approx 0$.
\end{itemize}
In small samples, both assumptions fail. We need the \alert{delta method}.
\end{frame}

% ============================================================
\section{Delta method}
% ============================================================

\begin{frame}{Delta method: the setup}
Suppose we have an estimator $\hat{\bm\theta} \in \mathbb{R}^k$ that is close to its target $\bm\theta$, with a known (or estimated) covariance matrix:
\[
\mathrm{Var}(\hat{\bm\theta}) = \Sigma.
\]
We care about a smooth scalar transformation $g(\bm\theta)$ --- e.g.\ a ratio, a log, a relative lift.

\medskip
\textbf{Question.} What is $\mathrm{Var}(g(\hat{\bm\theta}))$?

\medskip
\textbf{Problem.} In general, $g$ is nonlinear, so $\mathrm{Var}(g(\hat{\bm\theta}))$ has no closed form. But if $\hat{\bm\theta}$ is concentrated around $\bm\theta$, we can \alert{linearize} $g$ and get an approximation.
\end{frame}

\begin{frame}{Step 1: linearize $g$ via Taylor}
First-order Taylor expansion of $g$ around the true value $\bm\theta$:
\[
g(\hat{\bm\theta}) \;\approx\; g(\bm\theta) \;+\; \nabla g(\bm\theta)^\top (\hat{\bm\theta} - \bm\theta).
\]
Two important things to notice:
\begin{itemize}
  \item $g(\bm\theta)$ is a \alert{constant} (it's the truth, not random).
  \item $\nabla g(\bm\theta)$ is also a \alert{constant vector} (gradient evaluated at the truth).
  \item The only random object on the right-hand side is $\hat{\bm\theta} - \bm\theta$.
\end{itemize}
So we have reduced $g(\hat{\bm\theta})$ to (a constant) $+$ (a linear function of $\hat{\bm\theta}$).
\end{frame}

\begin{frame}{Step 2: variance of a linear function}
For any constant vector $\mathbf{a}$ and random vector $\mathbf{Z}$ with $\mathrm{Var}(\mathbf{Z}) = \Sigma$:
\[
\mathrm{Var}(\mathbf{a}^\top \mathbf{Z}) = \mathbf{a}^\top \Sigma \,\mathbf{a}.
\]
\textit{Quick derivation:}
\begin{align*}
\mathrm{Var}(\mathbf{a}^\top \mathbf{Z})
  &= \mathbb{E}\big[(\mathbf{a}^\top(\mathbf{Z}-\mu))^2\big] \\
  &= \mathbb{E}\big[\mathbf{a}^\top(\mathbf{Z}-\mu)(\mathbf{Z}-\mu)^\top \mathbf{a}\big] \\
  &= \mathbf{a}^\top \,\mathbb{E}[(\mathbf{Z}-\mu)(\mathbf{Z}-\mu)^\top]\, \mathbf{a}
   \;=\; \mathbf{a}^\top \Sigma\, \mathbf{a}.
\end{align*}
We will apply this with $\mathbf{a} = \nabla g(\bm\theta)$ and $\mathbf{Z} = \hat{\bm\theta}$.
\end{frame}

\begin{frame}{Step 3: put the two pieces together}
From Step 1:
\[
g(\hat{\bm\theta}) - g(\bm\theta) \;\approx\; \nabla g(\bm\theta)^\top (\hat{\bm\theta} - \bm\theta).
\]
Taking variance of both sides (the left-hand constant $g(\bm\theta)$ drops out):
\[
\mathrm{Var}\!\big(g(\hat{\bm\theta})\big)
\;\approx\; \mathrm{Var}\!\big(\nabla g(\bm\theta)^\top \hat{\bm\theta}\big).
\]
Applying Step 2 with $\mathbf{a} = \nabla g(\bm\theta)$:
\[
\boxed{\;\mathrm{Var}\!\left(g(\hat{\bm\theta})\right) \;\approx\; \nabla g(\bm\theta)^\top \,\Sigma\, \nabla g(\bm\theta).\;}
\]
In practice we plug in the estimates $\hat{\bm\theta}$ and $\hat\Sigma$ since $\bm\theta$ and $\Sigma$ are unknown.
\end{frame}

\begin{frame}{Delta method for a ratio}
Let $g(X, Y) = X / Y$. Then
\[
\nabla g = \begin{pmatrix} 1/Y \\ -X/Y^2 \end{pmatrix}.
\]
Plugging into $\nabla g^\top \Sigma \nabla g$:
\[
\mathrm{Var}\!\left(\frac{X}{Y}\right) \approx
\frac{\mathrm{Var}(X)}{Y^2}
+ \frac{X^2}{Y^4}\mathrm{Var}(Y)
- 2\,\frac{X}{Y^3}\mathrm{Cov}(X, Y).
\]
Here $X = \hat\tau$ (numerator) and $Y$ is the estimator of the control mean (denominator). What ``$Y$'' is concretely, and how we get $\mathrm{Var}(Y)$ and $\mathrm{Cov}(X, Y)$, depends on whether we have covariates.
\end{frame}

\begin{frame}{Recap}
We wanted a standard error for the relative lift
\[
\hat L \;=\; \frac{\hat\tau}{\bar Y_{\text{control}}}.
\]
The delta-method ratio formula tells us
\[
\mathrm{Var}(\hat L) \;\approx\;
\frac{\mathrm{Var}(\hat\tau)}{\bar Y_{\text{control}}^{\,2}}
+ \frac{\hat\tau^{\,2}}{\bar Y_{\text{control}}^{\,4}}\,\mathrm{Var}(\bar Y_{\text{control}})
- 2\,\frac{\hat\tau}{\bar Y_{\text{control}}^{\,3}}\,\mathrm{Cov}(\hat\tau,\bar Y_{\text{control}}).
\]
So computing $\mathrm{SE}(\hat L) = \sqrt{\mathrm{Var}(\hat L)}$ has been reduced to finding \alert{three} ingredients:
\begin{enumerate}
  \item $\mathrm{Var}(\hat\tau)$ \hfill --- variance of the absolute effect,
  \item $\mathrm{Var}(\bar Y_{\text{control}})$ \hfill --- variance of the control-mean estimator,
  \item $\mathrm{Cov}(\hat\tau, \bar Y_{\text{control}})$ \hfill --- their covariance.
\end{enumerate}
The rest of the talk is just: \alert{where do these three numbers come from?}
\end{frame}

\begin{frame}{Warning}


While simple to write, how can we easily calculate $\mathrm{Var}(\bar Y_{\text{control}})$ and $\mathrm{Cov}(\hat\tau, \bar Y_{\text{control}})$? We don't have a closed-form formula for any of them.

\end{frame}

% ============================================================
\section{Reminder: covariance matrix in OLS}
% ============================================================

\begin{frame}{Reminder: OLS hands us the full covariance matrix}
We are fitting OLS of the form $Y_i = \alpha + \tau T_i + \mathbf{X}_i^\top \bm\beta + \epsilon_i$.

Model coefficients are $\hat{\bm\theta} = (\hat\alpha,\, \hat\tau,\, \hat{\bm\beta}^\top)^\top$. Fitting OLS on the design matrix $X$ with residual variance $\hat\sigma^2$ gives:
\[
\widehat{\mathrm{Var}}(\hat{\bm\theta}) \;=\; \hat\sigma^2 (X^\top X)^{-1} \;=\; \Sigma_{\hat\theta}.
\]
This is a single $(2+p)\times(2+p)$ matrix containing:
\begin{itemize}
  \item diagonal: $\mathrm{Var}(\hat\alpha)$, $\mathrm{Var}(\hat\tau)$, $\mathrm{Var}(\hat\beta_j)$;
  \item off-diagonal: \emph{every} pairwise covariance, including
  $\mathrm{Cov}(\hat\tau,\hat\alpha)$ and $\mathrm{Cov}(\hat\tau,\hat\beta_j)$.
\end{itemize}
In \texttt{statsmodels}: \texttt{results.cov\_params()} returns exactly this matrix.

\medskip
As we'll see, both the variance of the denominator and its covariance with $\hat\tau$ are going to come out of $\Sigma_{\hat\theta}$.
\end{frame}

% ============================================================
\section{Two cases: no covariates vs. covariates}
% ============================================================

\begin{frame}{Case 1: no covariates}
Model: $Y_i = \alpha + \tau T_i + \epsilon_i$.
\begin{itemize}
  \item The control-group mean estimator is simply $\hat\alpha$: under $T=0$, $\mathbb{E}[Y\mid T=0] = \alpha$, and OLS estimates it with $\hat\alpha$.
  \item So in the ratio $\hat L = \hat\tau / \hat\alpha$, both $X = \hat\tau$ and $Y = \hat\alpha$ are \alert{OLS coefficients}.
  \item OLS already returns the full $2\times 2$ covariance matrix of $(\hat\alpha, \hat\tau)$:
  \[
  \widehat{\mathrm{Var}}\begin{pmatrix}\hat\alpha\\\hat\tau\end{pmatrix}
  = \hat\sigma^2 (X^\top X)^{-1}.
  \]
  \item So $\mathrm{Var}(\hat\tau)$, $\mathrm{Var}(\hat\alpha)$ and $\mathrm{Cov}(\hat\tau,\hat\alpha)$ are read off directly --- plug into the ratio formula and we are done.
\end{itemize}
\end{frame}

\begin{frame}{Case 2: with covariates --- the denominator is no longer a coefficient}
Model: $Y_i = \alpha + \tau T_i + \mathbf{X}_i^\top \bm\beta + \epsilon_i$.
\begin{itemize}
  \item Now $\hat\alpha$ alone is \alert{not} an estimate of the control mean. It is the predicted $Y$ when $\mathbf{X} = \mathbf{0}$ --- meaningless if covariates are not 0-centered.
  \item The natural estimator of the control mean is
  \[
  \bar{Y}^{\text{adj}}_{\text{control}} = \hat\alpha + \bar{\mathbf{X}}_{\text{control}}^\top \hat{\bm\beta},
  \]
  i.e.\ the model's prediction averaged over control covariates.
  \item This is a \alert{linear combination of several OLS coefficients}, not a single one.
  \item Consequences:
  \begin{itemize}
    \item $\mathrm{Var}(\bar{Y}^{\text{adj}}_{\text{control}})$ is not a single diagonal entry --- it is a quadratic form in $\Sigma_{\hat\theta}$.
    \item $\mathrm{Cov}(\hat\tau, \bar{Y}^{\text{adj}}_{\text{control}})$ is not a single off-diagonal entry --- it mixes $\mathrm{Cov}(\hat\tau,\hat\alpha)$ with all the $\mathrm{Cov}(\hat\tau, \hat\beta_j)$.
  \end{itemize}
  \item We need the machinery of variances of linear combinations.
\end{itemize}
\end{frame}

% ============================================================
\section{Linear model: variance of linear combinations}
% ============================================================

\begin{frame}{OLS gives us the full covariance matrix}
Stack parameters $\bm{\theta} = (\alpha, \tau, \bm{\beta}^\top)^\top$. OLS returns:
\[
\widehat{\mathrm{Var}}(\hat{\bm{\theta}}) = \hat\sigma^2 (X^\top X)^{-1}.
\]
We get \alert{every} variance \emph{and} every pairwise covariance ``for free'':
\[
\Sigma_{\hat\theta} =
\begin{pmatrix}
\mathrm{Var}(\hat\alpha)        & \mathrm{Cov}(\hat\alpha,\hat\tau)   & \mathrm{Cov}(\hat\alpha,\hat{\bm\beta}^\top) \\
\mathrm{Cov}(\hat\tau,\hat\alpha) & \mathrm{Var}(\hat\tau)            & \mathrm{Cov}(\hat\tau,\hat{\bm\beta}^\top)   \\
\mathrm{Cov}(\hat{\bm\beta},\hat\alpha) & \mathrm{Cov}(\hat{\bm\beta},\hat\tau) & \mathrm{Var}(\hat{\bm\beta})
\end{pmatrix}.
\]
This is the key object that makes a clean delta-method calculation possible.
\end{frame}

\begin{frame}{Variance of a linear combination of parameters}
For any constant vector $\mathbf{c}$:
\[
\mathrm{Var}(\mathbf{c}^\top \hat{\bm\theta}) = \mathbf{c}^\top \Sigma_{\hat\theta} \mathbf{c}.
\]
Expanding component-wise for $\mathbf{c}^\top \hat{\bm\theta} = \sum_j c_j \hat\theta_j$:
\[
\mathrm{Var}\!\left(\sum_j c_j \hat\theta_j\right)
= \sum_j c_j^2 \mathrm{Var}(\hat\theta_j) + 2\sum_{j<k} c_j c_k\,\mathrm{Cov}(\hat\theta_j, \hat\theta_k).
\]
And for the covariance with another linear combination $\mathbf{d}^\top \hat{\bm\theta}$:
\[
\mathrm{Cov}(\mathbf{c}^\top\hat{\bm\theta}, \mathbf{d}^\top\hat{\bm\theta}) = \mathbf{c}^\top \Sigma_{\hat\theta} \mathbf{d}.
\]
\end{frame}

% ============================================================
\section{Relative lift in OLS}
% ============================================================

\begin{frame}{The adjusted control mean}
With covariates, the predicted outcome for a control unit with covariates $\mathbf{x}$ is
\[
\hat{\mathbb{E}}[Y \mid T=0, \mathbf{X}=\mathbf{x}] = \hat\alpha + \mathbf{x}^\top \hat{\bm\beta}.
\]
Averaging over the empirical covariate distribution of the control group:
\[
\boxed{\;\bar{Y}^{\text{adj}}_{\text{control}} = \hat\alpha + \bar{\mathbf{X}}_{\text{control}}^\top \hat{\bm\beta}.\;}
\]
This is a linear combination $\mathbf{c}^\top \hat{\bm\theta}$ with
\[
\mathbf{c} = (1,\; 0,\; \bar{\mathbf{X}}_{\text{control}}^\top)^\top
\quad\text{(weight on } \hat\tau \text{ is } 0\text{).}
\]
\end{frame}

\begin{frame}{Each piece of the delta-method formula}
Same three ingredients we identified earlier, now written for the adjusted control mean $\bar{Y}^{\text{adj}}_{\text{control}}$ in the denominator:
\[
\mathrm{Var}(\hat L) \approx
\underbrace{\frac{\mathrm{Var}(\hat\tau)}{(\bar{Y}^{\text{adj}}_{\text{control}})^2}}_{\text{(A) Var}(\hat\tau)}
+ \underbrace{\frac{\hat\tau^2}{(\bar{Y}^{\text{adj}}_{\text{control}})^4}\,\mathrm{Var}(\bar{Y}^{\text{adj}}_{\text{control}})}_{\text{(B) Var}(\bar Y^{\text{adj}})}
- \underbrace{2\,\frac{\hat\tau}{(\bar{Y}^{\text{adj}}_{\text{control}})^3}\,\mathrm{Cov}(\hat\tau,\bar{Y}^{\text{adj}}_{\text{control}})}_{\text{(C) Cov}(\hat\tau,\bar Y^{\text{adj}})}.
\]
We now compute each of (A), (B), (C) as a function of entries of $\Sigma_{\hat\theta}$.
\end{frame}

\begin{frame}{(A) Variance of $\hat\tau$}
Read it off the OLS covariance matrix.
\[
\mathrm{Var}(\hat\tau) = \big[\Sigma_{\hat\theta}\big]_{\tau,\tau}.
\]
This is what \texttt{statsmodels} returns as \texttt{bse[treatment]**2}.
\end{frame}

\begin{frame}{(B) Variance of $\bar{Y}^{\text{adj}}_{\text{control}}$}
Quadratic form with $\mathbf{c} = (1, \bar{\mathbf{X}}_{\text{control}}^\top)^\top$ on the $(\alpha, \bm\beta)$ block:
\[
\mathrm{Var}(\bar{Y}^{\text{adj}}_{\text{control}}) =
\begin{pmatrix} 1 & \bar{\mathbf{X}}_{\text{control}}^\top \end{pmatrix}
\mathrm{Var}\!\begin{pmatrix} \hat\alpha \\ \hat{\bm\beta} \end{pmatrix}
\begin{pmatrix} 1 \\ \bar{\mathbf{X}}_{\text{control}} \end{pmatrix}.
\]
Expanded:
\[
\mathrm{Var}(\hat\alpha)
+ \bar{\mathbf{X}}_{\text{control}}^\top \mathrm{Var}(\hat{\bm\beta})\,\bar{\mathbf{X}}_{\text{control}}
+ 2\,\mathrm{Cov}(\hat\alpha, \hat{\bm\beta})^\top \bar{\mathbf{X}}_{\text{control}}.
\]
Every term lives in $\Sigma_{\hat\theta}$.
\end{frame}

\begin{frame}{(C) Covariance of $\hat\tau$ with the adjusted control mean}
$\hat\tau$ is linear in $\hat{\bm\theta}$ (with weight 1 on $\tau$, $0$ elsewhere); the control mean has weights $(1, \bar{\mathbf{X}}_{\text{control}})$ on $(\alpha, \bm\beta)$. So:
\[
\mathrm{Cov}(\hat\tau,\,\bar{Y}^{\text{adj}}_{\text{control}})
= \mathrm{Cov}(\hat\tau, \hat\alpha) + \sum_{j} \bar{X}_{\text{control},j}\,\mathrm{Cov}(\hat\tau, \hat\beta_j).
\]
\end{frame}

% ============================================================
\section{MDE for relative lift}
% ============================================================

\begin{frame}{What about relative MDE}
For an absolute effect, MDE at significance $\alpha$ and power $1-\beta$ is:
\[
\mathrm{MDE}_{\text{abs}} = (z_{1-\alpha/2} + z_{1-\beta}) \cdot \mathrm{SE}(\hat\tau).
\]
For relative lift one is tempted to write:
\[
\mathrm{MDE}_{\text{rel}} \stackrel{?}{=} (z_{1-\alpha/2} + z_{1-\beta}) \cdot \mathrm{SE}(\hat L).
\]
But $\mathrm{SE}(\hat L)$ \alert{itself depends on the true lift} --- this formula is circular. The correct way to solve for $\mathrm{MDE}_{\text{rel}}$ is to plug in the formula for $\mathrm{SE}(\hat L)$ and solve the resulting \alert{quadratic equation}. Luiz is working on this for the delta method, but looking for contributor for OLS.


This means that \alert{relative MDE sometimes does not exist}.
\end{frame}

% ============================================================
\section{In \texttt{cluster\_experiments}}
% ============================================================

\begin{frame}{All of this is wired up via \texttt{relative\_effect}}
Everything we derived (delta-method SE, adjusted control mean) is implemented behind a single flag:
\[
\texttt{relative\_effect: True}
\]
\begin{itemize}
  \item Available for \texttt{OLSAnalysis} and \texttt{ClusteredOLSAnalysis} (it needs the full $\Sigma_{\hat\theta}$).
  \item Switches the analyser to report $\hat L = \hat\tau / \bar{Y}^{\text{adj}}_{\text{control}}$ with delta-method SE instead of the raw $\hat\tau$.
  \item Carries through to \texttt{NormalPowerAnalysis}: \texttt{.mde(...)}\ returns the relative MDE, not the absolute one (not fully correct, does not solve the quadratic yet).
\end{itemize}
\end{frame}

\begin{frame}[fragile]{Config example 1: vanilla OLS, no covariates}
\small
\begin{verbatim}
config = {
    "analysis": "ols",
    "splitter": "non_clustered",
    "target_col": "orders",
    "relative_effect": True,
}

pw = NormalPowerAnalysis.from_dict(config)
pw.mde(user_df, n_simulations=1)   # -> relative MDE
\end{verbatim}
Here $Y = \hat\alpha$: Case 1 from earlier. Internally we use the $2\times 2$ covariance block of $(\hat\alpha, \hat\tau)$.
\end{frame}

\begin{frame}[fragile]{Config example 2: OLS with covariates}
\small
\begin{verbatim}
config = {
    "analysis": "ols",
    "splitter": "non_clustered",
    "target_col": "orders",
    "covariates": ["X1", "X2"],
    "relative_effect": True,
}

pw = NormalPowerAnalysis.from_dict(config)
pw.mde(user_df, n_simulations=1)
\end{verbatim}
Here $Y = \hat\alpha + \bar{\mathbf{X}}_{\text{control}}^\top \hat{\bm\beta}$: Case 2. The full covariance matrix of $(\hat\alpha, \hat\tau, \hat{\bm\beta})$ is used to build $\mathrm{Var}(Y)$ and $\mathrm{Cov}(\hat\tau, Y)$.
\end{frame}

\begin{frame}[fragile]{Config example 3: \texttt{AnalysisPlan} on real experiment data}
\small
Same flag, but on the inference side (analysing data from a finished test):
\begin{verbatim}
config = {
    "metrics": [
        {"alias": "Orders", "name": "orders"},
    ],
    "variants": [
        {"name": "A", "is_control": True},
        {"name": "B", "is_control": False},
    ],
    "analysis_type": "ols",
    "variant_col": "treatment",
    "analysis_config": {"relative_effect": True, "covariates": ["X1"]},
}

results = AnalysisPlan.from_metrics_dict(config).analyze(user_df)
\end{verbatim}
\texttt{results.ate} is the relative lift $\hat L$; \texttt{results.std\_error} is the delta-method SE; CIs and $p$-values follow.
\end{frame}

\begin{frame}{Live demo}
End-to-end runnable example (data simulation, absolute vs.\ relative SE, absolute vs.\ relative MDE):
\begin{center}
\href{https://david26694.github.io/cluster-experiments/relative.html\#experiment-analysis-with-absolute-and-relative-effects}{\alert{\texttt{cluster\_experiments} docs --- relative effects notebook}}
\end{center}
\small
\url{https://david26694.github.io/cluster-experiments/relative.html#experiment-analysis-with-absolute-and-relative-effects}
\end{frame}

\begin{frame}{Takeaways}
\begin{itemize}
  \item Relative lift is a \alert{ratio of estimators}; treat it as such.
  \item Delta method + full OLS covariance matrix gives the correct SE in one clean formula.
  \item For MDE estimation, solve the \alert{quadratic} --- linear scaling is only an approximation valid when the control mean is essentially noise-free.
  \item In \texttt{cluster\_experiments}, flip \texttt{relative\_effect: True} on \texttt{OLSAnalysis}/\texttt{ClusteredOLSAnalysis} --- it works for both \texttt{NormalPowerAnalysis} (planning, relative MDE) and \texttt{AnalysisPlan} (inference, relative ATE + SE).
\end{itemize}
\end{frame}

\end{document}
